\documentclass[10pt]{article}

\usepackage[preprint]{tmlr}
\usepackage{amsmath,amssymb,amsthm,mathtools}
\usepackage{aliascnt}
\usepackage{booktabs}
\usepackage{enumitem}
\usepackage{microtype}
\usepackage[hidelinks]{hyperref}
\usepackage{url}
\usepackage[nameinlink,noabbrev]{cleveref}

\newtheorem{definition}{Definition}[section]

\newaliascnt{theorem}{definition}
\newtheorem{theorem}[theorem]{Theorem}
\aliascntresetthe{theorem}

\newaliascnt{proposition}{definition}
\newtheorem{proposition}[proposition]{Proposition}
\aliascntresetthe{proposition}

\newaliascnt{lemma}{definition}
\newtheorem{lemma}[lemma]{Lemma}
\aliascntresetthe{lemma}

\newaliascnt{corollary}{definition}
\newtheorem{corollary}[corollary]{Corollary}
\aliascntresetthe{corollary}

\newaliascnt{example}{definition}
\newtheorem{example}[example]{Example}
\aliascntresetthe{example}

\newaliascnt{remark}{definition}
\newtheorem{remark}[remark]{Remark}
\aliascntresetthe{remark}

\crefname{definition}{definition}{definitions}
\crefname{theorem}{theorem}{theorems}
\crefname{proposition}{proposition}{propositions}
\crefname{lemma}{lemma}{lemmas}
\crefname{corollary}{corollary}{corollaries}
\crefname{example}{example}{examples}
\crefname{remark}{remark}{remarks}
\Crefname{definition}{Definition}{Definitions}
\Crefname{theorem}{Theorem}{Theorems}
\Crefname{proposition}{Proposition}{Propositions}
\Crefname{lemma}{Lemma}{Lemmas}
\Crefname{corollary}{Corollary}{Corollaries}
\Crefname{example}{Example}{Examples}
\Crefname{remark}{Remark}{Remarks}

\newcommand{\Id}{\operatorname{id}}
\newcommand{\Ker}{\operatorname{ker}}
\newcommand{\MinSupp}{\operatorname{MinSupp}}
\newcommand{\Proc}{\mathfrak{P}}
\newcommand{\im}{\operatorname{im}}
\newcommand{\Rcal}{\mathcal{R}}
\newcommand{\Ccal}{\mathcal{C}}

\title{Mechanistic Localization Is Support-Relative}

\author{
\name Luis F. Rosario Freytes\thanks{This research was conducted independently; the University of Michigan affiliation is provided solely to identify the author's current institution.}\\
\addr University of Michigan\\
\addr Ann Arbor, MI 48109, USA\\
\email luisrosa@umich.edu
}

\begin{document}
\maketitle

\begin{abstract}
Mechanistic-localization studies compare internal components across prompt classes, datasets, ablation schemes, and intervention protocols, but the domain on which a localization claim is evaluated is often treated as experimental metadata rather than as an argument of the claim. We make this domain explicit as an analysis support: source regimes generate compatible supports across represented cuts, while interventions may specify supports directly at an internal cut. For a fixed support $S$, phenomenon $\Phi:S\to Y_\Phi$, and internal access $L:S\twoheadrightarrow U$, exact realization holds if and only if $\Ker L\subseteq\Ker\Phi$. Localization therefore depends on which supported state pairs the access must distinguish. Restricting support can only enlarge the feasible realization family, so localization order can decrease and minimal realization identity can change without any change to the weights. Across nonnested supports, source overlap can understate internal-state overlap, and separate local realizations extend to the union exactly when cross-support access collisions are phenomenon-consistent. These results make support an explicit argument of mechanistic-localization claims and give exact checks for comparing localizations across contexts.\footnotemark
\end{abstract}

\footnotetext{AI tools were used to assist with drafting and editing. The author is responsible for the ideas, claims, proofs, citations, and final text.}

\section{Introduction}

A mechanistic-localization claim is never determined by the network parameters alone. A study also fixes a behavior to explain, a collection of natural or counterfactual states on which the claim is evaluated, and a decomposition of the network into candidate internal accesses. Recent work makes the dependence on these choices visible through ablation-sensitive faithfulness, input-conditioned circuit variation, explicit verification domains, and intervention-induced hidden states \citep{miller2024faithfulness,bayat2026many,hadad2026formal,makelov2024subspace,grant2026divergent}.

These results are usually discussed as problems of faithfulness, identifiability, minimality, transfer, or intervention design. They also expose a common mathematical dependence. A trained network defines maps on ambient state spaces. A localization analysis evaluates those maps only on a selected set of states. Changing that set changes which state pairs can be compared and which relations among coordinates hold on the domain. An internal access that is insufficient on the ambient state space can therefore become sufficient after restriction even though the parameters and ambient maps are unchanged.

We take the evaluation domain to be an argument of the localization claim. Fix a deterministic network with parameter setting $\theta$. A source regime $C\subseteq X_0$ generates a reachable support
\[
S_q^{\theta,C}=A_q^\theta(C)
\]
at cut $q$, where $A_q^\theta$ is the network prefix. Intervention protocols can instead supply a set of states directly at an internal cut. The theory below distinguishes these source-generated supports from general analysis supports. When one source regime is propagated through all represented cuts, the resulting restricted states, receiver exposures, and continuations form a \emph{supported neural process}.

The support determines which internal distinctions are available for comparison, but it does not determine which of them matter for the behavior being explained. We therefore also fix a phenomenon profile
\[
\Phi:S\to Y_\Phi
\]
on the analysis support $S$ and a declared internal access
\[
L:S\twoheadrightarrow U.
\]
Exact realization means that the phenomenon can be recovered from the accessed state:
\begin{equation}
\Phi=hL
\quad\Longleftrightarrow\quad
\Ker L\subseteq\Ker\Phi.
\label{eq:intro-realization}
\end{equation}
The two kernels are partitions of the same supported domain. The access may identify states only when the declared phenomenon identifies them as well.

Once localization is stated on this object, several forms of context dependence follow from ordinary restriction. If $S\subseteq S'$, an access that realizes the phenomenon on $S'$ also realizes its restriction to $S$. For receiver families, the feasible localization set can therefore grow as support narrows. The minimum receiver count can decrease, a different family can become minimal, and several minimal realizations can appear where only one existed on the larger domain. A dependency can also disappear when the phenomenon becomes constant on the restricted support.

Nonnested contexts expose two further distinctions. Input overlap need not equal internal-state overlap because a noninjective network prefix can map different source inputs to the same hidden state. Conversely, a localization that works separately on two supports need not work on their union: the union introduces cross-support state pairs that neither local claim had to distinguish. Both effects concern the domain on which the mechanistic claim is stated, not a change in the ambient network.

The nearby literature varies several of these ingredients together. \Cref{sec:related-work} separates those axes; the formal development then varies support while holding the network, declared phenomenon, and internal access family fixed.

The analysis is exact and set-theoretic. Represented cuts have finitely many declared coordinates, but coordinate values may be arbitrary, including ordinary real-valued activations. We do not assume linearity, differentiability, a finite activation alphabet, or a probability distribution on the support. The access decomposition is declared by the analysis. Causal necessity, algorithmic abstraction, and equivalence between different presentation languages require additional structure beyond the exact interface-realization criterion studied here.

\section{Localization depends on an evaluation domain}
\label{sec:related-work}

The paper isolates one variable that is present across several parts of mechanistic interpretability: the set of states on which a localization claim is evaluated. Existing studies vary the evaluation domain as well as other ingredients such as ablation rules, mediator choices, prompt statistics, or verification conditions. We separate those ingredients here so that the theory can later hold the network, phenomenon, and declared access family fixed and vary only support.

\subsection{Circuit variation across inputs and interventions}

Circuit extraction already shows that localization depends on how the model is interrogated. \citet{miller2024faithfulness} show that circuit-faithfulness scores can change substantially with ablation details. \citet{bayat2026many} vary input-token statistics while holding the task fixed and recover structurally different circuits that nevertheless transfer across conditions and implement the same computation. These results do not by themselves isolate support as the cause of the observed variation, but they show that a circuit claim is evaluated relative to a chosen input regime and evaluation procedure.

Interventions make the domain dependence more explicit. \citet{makelov2024subspace} construct cases in which a subspace intervention obtains the desired end-to-end effect by activating a dormant parallel pathway. \citet{grant2026divergent} analyze intervention-induced representations that depart from the natural representation distribution, distinguishing cases that are harmless for specified functional claims from cases that activate dormant behavioral pathways. Verification-based circuit discovery likewise states guarantees on explicit continuous input and patching domains \citep{hadad2026formal}. In each case, the set of internal states admitted by the analysis is part of the condition under which the localization claim is evaluated.

\subsection{Localization criteria, mediator choice, and non-uniqueness}

The evaluation domain is only one argument of a mechanistic claim. Other work studies what should count as a candidate circuit, mediator, or high-level mechanism and what criterion should certify it. Causal abstraction formalizes relations between lower-level causal systems and higher-level mechanistic models \citep{geiger2025causal}. The causal-mediation perspective of \citet{mueller2026mediator} emphasizes that mechanistic studies choose different causal units or mediators. MIB standardizes circuit and causal-variable localization across tasks and models \citep{mueller2025mib}, while \citet{shi2024hypothesis} formulate hypothesis tests for proposed circuits and \citet{adolfi2025complexity} analyze the complexity of circuit-discovery queries.

Minimality and identity introduce further choices. \citet{hadad2026formal} distinguish several notions of minimality within a verification-based setting. \citet{meloux2025identifiable} exhibit multiple circuits and mechanistic interpretations satisfying the same declared validity criteria in small models. \citet{sun2026equivalent} formalize equivalence between mechanistic interpretations across neural networks. These questions remain after a support has been fixed. Conversely, support dependence remains after a localization criterion and access decomposition have been fixed.

A closely related formalism studies neural architecture through represented receiver processes $B_j=G_jQ_j$, using interface fibers to characterize extensional distinction structure, unmarked factorization identity, marked re-presentation, continuation accessibility, and architectural dependence on upstream composition \citep{rosariofreytes2026architecture}. The present paper asks a different question on the same general interface substrate. We fix the ambient network, phenomenon, and declared access family and vary the analysis support, then study the resulting family of realizing receiver sets and its behavior under restriction and combination of supports.

\subsection{The missing argument}

Across these settings, the network, phenomenon, localization procedure, mediator family, or intervention protocol may all be stated explicitly. The supported internal domain is often induced or specified operationally by a dataset, prompt family, corruption set, patching range, or verification region, but it is not always included in the mathematical type of the localization claim. Verification work can certify a proposed circuit on a specified domain; the present question is how localization itself changes when the domain changes while the ambient network, phenomenon, and declared access family are held fixed.

The theory below isolates that argument. We fix the ambient network and a declared family of internal accesses, fix the phenomenon whose distinctions must be preserved, and ask how the feasible localization structure changes when only the analysis support changes. This turns several observations that otherwise appear under different headings into consequences of one operation: restriction of the state pairs on which exact realization is tested. \Cref{thm:realization-monotonicity} gives the nested-support direction, while \cref{sec:context-comparison} treats overlap and union questions for nonnested contexts.

\section{The object of localization: a supported neural process}
\label{sec:supported-process}

A fixed parameter tensor does not determine a unique localization problem. The network specifies an ambient computation, but an analysis supplies a domain on which internal distinctions are compared. This section makes that separation explicit before introducing the realization criterion.

\subsection{A continuous example}

Consider a represented cut with state
\[
x=(x_1,x_2)\in\mathbb R^2
\]
and downstream behavior
\begin{equation}
\Phi(x_1,x_2)=x_1+x_2.
\label{eq:continuous-preview}
\end{equation}
This is the computation of a single linear neuron with weights $(1,1)$. Take the two coordinate accesses
\[
L_1(x)=x_1,
\qquad
L_2(x)=x_2.
\]
When discussing receiver-family minimality below, we use these two coordinate accesses as the declared two-receiver decomposition. On the full ambient domain $\mathbb R^2$, neither coordinate alone determines $\Phi$: holding $x_1$ fixed while varying $x_2$ changes the output, and conversely. The joint access $(L_1,L_2)$ is required.

Now restrict the same neuron to the diagonal
\[
D=\{(t,t):t\in\mathbb R\}.
\]
On $D$,
\[
\Phi(t,t)=2t=2L_1(t,t)=2L_2(t,t).
\]
Either coordinate alone now determines the behavior. On the vertical line
\[
V=\{(0,t):t\in\mathbb R\},
\]
only $L_2$ is needed. The weights, ambient state space, and formula \eqref{eq:continuous-preview} have not changed. The localization structure changes because the supported states satisfy different relations.

A localization statement compares the distinctions preserved by an internal access with the distinctions that remain possible on the domain of analysis. The relevant object is therefore the ambient neural map together with the support on which it is interrogated.

\subsection{Ambient neural computation}

Fix a deterministic feed-forward computation with represented state cuts
\[
q=0,1,\ldots,L.
\]
At cut $q$, let $I_q$ be a finite coordinate index set. Coordinate $i\in I_q$ takes values in a nonempty set $A_{q,i}$, and
\[
X_q:=\prod_{i\in I_q}A_{q,i}.
\]
No finiteness assumption is placed on $A_{q,i}$.

A fixed trained parameter setting is denoted by $\theta$. For $q<L$, the represented layer update is
\[
F_q^\theta:X_q\to X_{q+1}.
\]
Write
\[
A_0^\theta=\Id_{X_0},
\qquad
A_q^\theta=F_{q-1}^\theta\circ\cdots\circ F_0^\theta
\quad(q>0)
\]
for the network prefix to cut $q$.

Localization also requires a declared decomposition into candidate internal units. For every receiver $j\in I_{q+1}$, let
\[
F_{q,j}^\theta:=\pi_{q+1,j}F_q^\theta:X_q\to A_{q+1,j}.
\]
We assume that the implementation under study supplies an exposure map and continuation
\begin{equation}
X_q\xrightarrow{Q_{q,j}^\theta}Z_{q,j}
\xrightarrow{G_{q,j}^\theta}A_{q+1,j},
\qquad
F_{q,j}^\theta=G_{q,j}^\theta Q_{q,j}^\theta.
\label{eq:receiver-factorization}
\end{equation}
The exposure $Q_{q,j}^\theta$ is not an arbitrary factorization inserted after the fact. It names the state available at the receiver boundary chosen by the analysis. This interface/continuation grain is the represented receiver-process object used for architectural individuation in \citet{rosariofreytes2026architecture}. Here it provides the ambient access structure that will be restricted to an analysis support and tested against a declared phenomenon. A neuron-, head-, edge-, feature-, or mediator-level study can therefore induce different access maps even for the same trained network.

\subsection{Reachable support}

\begin{definition}[Source regime and reachable support]
For any subset $C\subseteq X_0$, define its reachable support at cut $q$ by
\begin{equation}
S_q^{\theta,C}:=A_q^\theta(C)\subseteq X_q.
\label{eq:reachable-support}
\end{equation}
A \emph{source regime} is a nonempty such subset. We allow $C=\varnothing$ only when forming intersections of regimes; then $S_q^{\theta,\varnothing}=\varnothing$.
\end{definition}

\begin{definition}[Analysis support at a cut]
\label{def:analysis-support}
Fix a cut $q$. An \emph{analysis support} is a nonempty subset
\[
S\subseteq X_q
\]
on which a localization claim is evaluated. It is \emph{source-generated} when $S=S_q^{\theta,C}$ for some source regime $C$. An intervention protocol may instead specify states directly at cut $q$, in which case the resulting analysis support need not be the prefix image of any source regime.
\end{definition}

At a fixed cut, an ambient receiver exposure, internal access, or phenomenon can be restricted directly to any analysis support $S$. The factorization and kernel criteria below depend only on that restricted domain. Statements that compare several cuts require a compatible family of supports; a source-generated family $\{S_q^{\theta,C}\}_{q=0}^L$ supplies such compatibility automatically through the network updates.

The ambient map $F_q^\theta$ is defined on all of $X_q$. Source-generated supports satisfy
\[
F_q^\theta(S_q^{\theta,C})=S_{q+1}^{\theta,C}.
\]
We therefore restrict and corestrict the update to
\[
F_q^{\theta,C}:
S_q^{\theta,C}\to S_{q+1}^{\theta,C},
\qquad
F_q^{\theta,C}(x):=F_q^\theta(x).
\]
Its receiver component is
\[
F_{q,j}^{\theta,C}
:=\pi_{q+1,j}F_q^{\theta,C}
=F_{q,j}^\theta|_{S_q^{\theta,C}}.
\]
For $q\le r$, write
\begin{equation}
F_{q\rightsquigarrow r}^{\theta,C}
:=
\begin{cases}
\Id_{S_q^{\theta,C}}, & q=r,\\
F_{r-1}^{\theta,C}\circ\cdots\circ F_q^{\theta,C}, & q<r,
\end{cases}
\label{eq:supported-continuation}
\end{equation}
for the supported continuation from cut $q$ to cut $r$.
If $C\subseteq C'$, then
\begin{equation}
S_q^{\theta,C}\subseteq S_q^{\theta,C'}
\qquad\text{for every }q.
\label{eq:support-monotonicity-basic}
\end{equation}
The inclusion is elementary, but it is the source of the later monotonicity theorem: a smaller regime creates fewer pairs of internal states that an access must distinguish.

For receiver $j$, define the realized image
\[
Z_{q,j}^{\theta,C}:=
Q_{q,j}^\theta(S_q^{\theta,C})
\]
and corestrict the exposure to
\[
Q_{q,j}^{\theta,C}:
S_q^{\theta,C}\twoheadrightarrow Z_{q,j}^{\theta,C}.
\]
Restricting $G_{q,j}^\theta$ to this realized image, define
\[
G_{q,j}^{\theta,C}
:=
G_{q,j}^\theta\big|_{Z_{q,j}^{\theta,C}}
:
Z_{q,j}^{\theta,C}\to A_{q+1,j}.
\]
Then
\[
F_{q,j}^{\theta,C}
=
G_{q,j}^{\theta,C}Q_{q,j}^{\theta,C}.
\]

\begin{definition}[Supported neural process]
For fixed $\theta$ and $C$, the \emph{supported neural process} is
\begin{equation}
\Proc_{\mathrm{NN}}(\theta,C)
:=
\left(
\{S_q^{\theta,C}\}_{q=0}^L,
\{Q_{q,j}^{\theta,C}\}_{q<L,\,j\in I_{q+1}},
\{G_{q,j}^{\theta,C}\}_{q<L,\,j\in I_{q+1}}
\right).
\label{eq:supported-process}
\end{equation}
\end{definition}

The weights determine the ambient maps. The regime determines which parts of their domains are realized. The supported process retains both pieces explicitly, so changing the regime changes the domain of the internal maps without changing their ambient definitions.

\subsection{Receiver fibers}

For a map $f:S\to T$, write
\[
\Ker f
:=
\{(x,x')\in S\times S:f(x)=f(x')\}.
\]
This is the equivalence relation induced by the fibers of $f$. For an internal access, the fibers are the distinctions it removes: two supported states in the same fiber are presented identically downstream of that access.

At receiver $(q,j)$ define
\begin{equation}
D_{q,j}^{\theta,C}:=
\Ker Q_{q,j}^{\theta,C}.
\label{eq:receiver-kernel}
\end{equation}
Because $Q_{q,j}^{\theta,C}$ is a restricted map, its kernel is a relation on the supported domain $S_q^{\theta,C}\times S_q^{\theta,C}$; changing the support changes the state pairs on which the relation is defined even when the ambient exposure $Q_{q,j}^{\theta}$ is unchanged.

The quotient
\[
S_q^{\theta,C}/D_{q,j}^{\theta,C}
\]
is isomorphic to the realized receiver image $Z_{q,j}^{\theta,C}$. The quotient notation will be useful because localization depends on which states an access identifies, not on a preferred coordinate system for its codomain.

From \eqref{eq:receiver-factorization},
\begin{equation}
\Ker Q_{q,j}^{\theta,C}
\subseteq
\Ker F_{q,j}^{\theta,C}.
\label{eq:receiver-no-recovery}
\end{equation}
A distinction absent from the exposed receiver state cannot affect that receiver's immediate output through the fixed continuation. The rest of the paper asks which of these retained distinctions are needed for a declared phenomenon.

\section{Support-relative dependency}
\label{sec:support-dependency}

The continuous example in \cref{sec:supported-process} already shows that coordinate dependence is not determined by an ambient formula alone. This section records the exact criterion for coordinate sufficiency and one consequence for circuit-style necessity tests.

For $K\subseteq I_q$, let
\[
X_{q,K}:=\prod_{i\in K}A_{q,i}
\]
and let $\pi_K:X_q\to X_{q,K}$ be coordinate projection. For $K=\varnothing$, take $X_{q,\varnothing}$ to be a singleton.

\begin{definition}[Sufficient coordinate set]
A subset $K\subseteq I_q$ is \emph{sufficient for receiver $j$ on $C$} if there exists
\[
\widehat Q_{q,j,K}^{\theta,C}:
\pi_K(S_q^{\theta,C})\to Z_{q,j}^{\theta,C}
\]
such that
\begin{equation}
Q_{q,j}^{\theta,C}
=
\widehat Q_{q,j,K}^{\theta,C}
\,\pi_K|_{S_q^{\theta,C}}.
\label{eq:coordinate-sufficiency}
\end{equation}
\end{definition}

The definition asks whether the receiver state can be reconstructed from the coordinates in $K$ on the supported domain. It does not ask whether the ambient expression for $Q_{q,j}^\theta$ contains those coordinates syntactically. Relations that hold only on the support can make one coordinate determine another.

\begin{proposition}[Coordinate sufficiency criterion]
\label{prop:coordinate-sufficiency}
A coordinate set $K$ is sufficient for receiver $j$ on $C$ if and only if
\begin{equation}
\Ker(\pi_K|_{S_q^{\theta,C}})
\subseteq
\Ker Q_{q,j}^{\theta,C}.
\label{eq:coordinate-kernel-criterion}
\end{equation}
\end{proposition}

\begin{proof}
If \eqref{eq:coordinate-sufficiency} holds, equal projected values produce equal receiver values. Conversely, suppose \eqref{eq:coordinate-kernel-criterion} holds. Define
\[
\widehat Q_{q,j,K}^{\theta,C}(\pi_K(x))
:=Q_{q,j}^\theta(x).
\]
The kernel inclusion makes this independent of the representative $x\in S_q^{\theta,C}$.
\end{proof}

Because $I_q$ is finite and $I_q$ itself is sufficient, every receiver has at least one inclusion-minimal sufficient coordinate set. Write
\[
\MinSupp_{q,j}(\theta,C)
\]
for the collection of these minimal sets.

Minimality inherits the domain dependence of sufficiency. On the diagonal in the continuous example, either coordinate is a minimal sufficient access for the linear readout even though both are required on $\mathbb R^2$. Restriction can also remove a dependency entirely. For
\[
Q(a,b)=a\oplus b
\]
on $\{0,1\}^2$, both coordinates are needed on the full square, while $Q$ is constant on the diagonal $\{(0,0),(1,1)\}$. The empty coordinate set is therefore sufficient after restriction.

More generally, sufficiency is monotone under restriction. If $S\subseteq S'\subseteq X_q$ and $K$ is sufficient for a fixed map $Q$ on $S'$, then $K$ is sufficient on $S$. The same kernel inclusion is simply being tested on fewer pairs. Minimality need not be preserved, because a proper subset of $K$ may become sufficient on the smaller domain.

\subsection{One-coordinate operativity and necessity}

A common local test varies one component while holding the others fixed and asks whether a downstream quantity changes. On a restricted support, such a comparison is available only when both states belong to the support.

\begin{definition}[Direct supported operativity]
For coordinate $i\in I_q$ and receiver $j$, write
\[
i\to_{\theta,C}j
\]
when there exist $x,x'\in S_q^{\theta,C}$ such that
\[
x_k=x'_k\quad\text{for all }k\neq i,
\qquad
Q_{q,j}^\theta(x)\neq Q_{q,j}^\theta(x').
\]
\end{definition}

\begin{proposition}[Direct operativity and minimal sufficient supports]
\label{prop:direct-necessity}
For fixed $(q,j,\theta,C)$ and $i\in I_q$,
\begin{equation}
i\to_{\theta,C}j
\quad\Longleftrightarrow\quad
i\in K
\text{ for every }
K\in\MinSupp_{q,j}(\theta,C).
\label{eq:direct-necessity}
\end{equation}
\end{proposition}

\begin{proof}
Suppose $i\to_{\theta,C}j$, witnessed by $x,x'$. Any coordinate set $K$ omitting $i$ gives $\pi_K(x)=\pi_K(x')$ while the receiver values differ, so $K$ is not sufficient. Every minimal sufficient set must therefore contain $i$.

Conversely, suppose no direct witness exists. Whenever two supported states agree on all coordinates except possibly $i$, their receiver values agree. Hence $I_q\setminus\{i\}$ is sufficient by \cref{prop:coordinate-sufficiency}. Since $I_q$ is finite, it contains an inclusion-minimal sufficient subset omitting $i$.
\end{proof}

The proposition separates several circuit claims that are easy to conflate. A one-coordinate witness characterizes membership in every inclusion-minimal sufficient coordinate set. Membership in at least one minimal set, attribution thresholds, and interventions performed on a larger counterfactual domain are different criteria.

\section{Phenomenon-relative realization}
\label{sec:phenomenon-realization}

Coordinate sufficiency concerns what is needed to reproduce a receiver state. Mechanistic localization usually asks a different question: which internal state is sufficient for a behavior or other phenomenon of interest? An access can preserve many distinctions that the phenomenon ignores. The phenomenon must therefore be part of the localization claim.

Fix a regime $C$ and a stage $r$. A \emph{phenomenon profile} is a map
\[
\Phi_r^{\theta,C}:S_r^{\theta,C}\to Y_\Phi.
\]
The codomain may represent a scalar output, a vector of observables, a class label, or a tuple of tests. For $q\le r$, pull the profile back through the supported continuation:
\begin{equation}
\Phi_q^{\theta,C}
:=
\Phi_r^{\theta,C}
F_{q\rightsquigarrow r}^{\theta,C}.
\label{eq:phenomenon-pullback}
\end{equation}
The map $\Phi_q^{\theta,C}$ records which distinctions at cut $q$ matter for the declared downstream profile.

Its fibers define the corresponding phenomenon content. Let
\[
D_{\Phi,q}^{\theta,C}:=
\Ker\Phi_q^{\theta,C},
\qquad
\Ccal_{\Phi,q}^{\theta,C}
:=S_q^{\theta,C}/D_{\Phi,q}^{\theta,C}.
\]
Two supported states lie in the same class exactly when the declared phenomenon treats them identically. The quotient is therefore the coarsest partition of the supported state that still determines the phenomenon. It does not localize that content: it is defined using the whole state at cut $q$.

Now let
\[
L:S_q^{\theta,C}\twoheadrightarrow U
\]
be a declared internal access, corestricted to its realized image. Coordinate projections, individual receiver exposures, joint receiver states, and feature-level maps are all examples. Different choices of $L$ ask different localization questions even when the network, support, and phenomenon are fixed.

\begin{definition}[Exact interface realization]
The access $L$ \emph{realizes $\Phi$ on $C$} if there exists a map
\[
\widehat\Phi_L:U\to Y_\Phi
\]
such that
\begin{equation}
\Phi_q^{\theta,C}=\widehat\Phi_L L.
\label{eq:realization-definition}
\end{equation}
\end{definition}

The definition says that once the accessed state is known, no additional distinction in the supported cut state is needed to determine the phenomenon. The criterion can be stated entirely in terms of fibers.

This is the standard factorization-through-fibers criterion specialized to a phenomenon and internal access on a common supported domain. The same criterion is used for receiver-interface factorization in \citet{rosariofreytes2026architecture}: a map $R$ factors through a surjective interface $Q$ exactly when $\Ker Q\subseteq\Ker R$.

\begin{theorem}[Interface realization criterion]
\label{thm:realization-criterion}
An access
\[
L:S_q^{\theta,C}\twoheadrightarrow U
\]
realizes $\Phi$ if and only if
\begin{equation}
\boxed{
\Ker L
\subseteq
\Ker\Phi_q^{\theta,C}.
}
\label{eq:realization-kernel}
\end{equation}
\end{theorem}

\begin{proof}
If \eqref{eq:realization-definition} holds, two states in the same fiber of $L$ have the same phenomenon value. Hence $\Ker L\subseteq\Ker\Phi_q^{\theta,C}$. Conversely, suppose the kernel inclusion holds. Define
\[
\widehat\Phi_L(L(x)):=\Phi_q^{\theta,C}(x).
\]
The inclusion makes this independent of the representative $x$, and surjectivity of $L$ defines $\widehat\Phi_L$ on all of $U$.
\end{proof}

The proof uses only the common domain of $L$ and $\Phi$. Hence the same criterion applies at any fixed-cut analysis support $S\subseteq X_q$ from \cref{def:analysis-support}: for maps $L:S\twoheadrightarrow U$ and $\Phi:S\to Y_\Phi$, exact realization is equivalent to $\Ker L\subseteq\Ker\Phi$.

On a fixed support, the criterion settles exact interface realization. The support-relative problem begins when the same ambient access and phenomenon are restricted to different domains. The next sections construct the resulting realizing families and study their minimality, monotonicity under restriction, and compatibility across nonnested supports.

The access $L$ partitions the supported state according to what the proposed internal interface preserves. The phenomenon partitions the same state according to what must remain distinguishable. Realization holds exactly when the access never collapses a distinction that the phenomenon requires.

For an implemented receiver $j$ this becomes
\begin{equation}
Q_{q,j}^{\theta,C}\text{ realizes }\Phi
\quad\Longleftrightarrow\quad
D_{q,j}^{\theta,C}
\subseteq
D_{\Phi,q}^{\theta,C}.
\label{eq:receiver-realization}
\end{equation}
A phenomenon can therefore be determined by the whole layer state while failing to factor through the state exposed to a selected receiver.

This is an exact interface-sufficiency statement. It does not by itself establish causal necessity, a unique circuit, or a human-readable explanation. Those claims require additional intervention or representation structure. Here the factorization criterion isolates the lower-level condition that the proposed access retain the distinctions needed for the declared phenomenon.

\section{Distributed and minimal localization}
\label{sec:distributed}

A phenomenon need not be available at any single receiver. Once realization is phrased as a factorization problem, the same criterion applies to collections of receivers without introducing an attribution score or a new notion of sufficiency.

For $K\subseteq I_{q+1}$, define the joint exposure
\[
Q_{q,K}^{\theta,C}
:=
\langle Q_{q,j}^{\theta,C}\rangle_{j\in K}
\]
corestricted to its realized image. For $K=\varnothing$, take the unique map to a singleton. Two supported states have the same joint exposure exactly when every receiver in $K$ sees them as the same state, so
\begin{equation}
\Ker Q_{q,K}^{\theta,C}
=
\bigcap_{j\in K}D_{q,j}^{\theta,C},
\label{eq:joint-kernel}
\end{equation}
where the empty intersection is the universal relation on $S_q^{\theta,C}$.

Combining \eqref{eq:joint-kernel} with \cref{thm:realization-criterion}, the receiver family $K$ realizes $\Phi$ exactly when
\begin{equation}
\bigcap_{j\in K}D_{q,j}^{\theta,C}
\subseteq
D_{\Phi,q}^{\theta,C}.
\label{eq:family-realization}
\end{equation}
This expression makes the role of a distributed family explicit. Each receiver removes some distinctions. The family succeeds when their joint state separates every pair that the phenomenon needs separated.

The joint next-layer exposure factors the complete update. Let
\[
Q_q^{\theta,C}:=Q_{q,I_{q+1}}^{\theta,C}.
\]
On its realized joint image, the receiver continuations assemble into a map $G_q^{\theta,C}$ with
\begin{equation}
F_q^{\theta,C}=G_q^{\theta,C}Q_q^{\theta,C}.
\label{eq:joint-factorization}
\end{equation}
If $q<r$ and $\Phi_q^{\theta,C}$ is the pullback of a profile at stage $r$, then
\[
\Phi_q^{\theta,C}
=
\Phi_r^{\theta,C}
F_{q+1\rightsquigarrow r}^{\theta,C}
G_q^{\theta,C}Q_q^{\theta,C}.
\]
Hence
\begin{equation}
\Ker Q_q^{\theta,C}
\subseteq
\Ker\Phi_q^{\theta,C}.
\label{eq:full-joint-realization}
\end{equation}
The declared next-layer receiver organization jointly retains every distinction required for a later phenomenon, although no single receiver need do so.

\begin{definition}[Distributed receiver realization]
A phenomenon is \emph{distributed at cut $q$ relative to the declared receiver decomposition} if no individual receiver realizes it but some family $K$ with $|K|>1$ satisfies \eqref{eq:family-realization}.
\end{definition}

For example, on $S=\{0,1\}^2$ with
\[
\Phi(a,b)=a\oplus b,
\qquad
Q_1(a,b)=a,
\qquad
Q_2(a,b)=b,
\]
neither receiver realizes parity individually, while the joint exposure $(a,b)$ does.

\subsection{Realizing families and localization order}

Define the collection of realizing receiver families by
\begin{equation}
\Rcal_{\Phi,q}^{\theta,C}
:=
\left\{
K\subseteq I_{q+1}:
\bigcap_{j\in K}D_{q,j}^{\theta,C}
\subseteq
D_{\Phi,q}^{\theta,C}
\right\}.
\label{eq:realizing-families}
\end{equation}
It is upward closed: once $K$ realizes the phenomenon, adding receivers can only refine the joint fibers.

The minimum receiver count is
\begin{equation}
d_{\Phi,q}^{\theta,C}
:=
\begin{cases}
\min\{|K|:K\in\Rcal_{\Phi,q}^{\theta,C}\},
&\Rcal_{\Phi,q}^{\theta,C}\neq\varnothing,\\[0.3em]
+\infty,&\text{otherwise.}
\end{cases}
\label{eq:realization-order}
\end{equation}
For a nonconstant downstream profile with $q<r$, \eqref{eq:full-joint-realization} gives
\[
1\le d_{\Phi,q}^{\theta,C}\le |I_{q+1}|.
\]
A constant profile has order $0$.

The same condition has an equivalent pair-cover form. Define the phenomenon-distinguished pairs
\[
P_{\Phi,q}^{\theta,C}
:=
\left\{
(x,x')\in (S_q^{\theta,C})^2:
\Phi_q^{\theta,C}(x)\neq\Phi_q^{\theta,C}(x')
\right\},
\]
and for each receiver $j$ let
\[
E_{q,j}^{\theta,C}
:=
\left\{
(x,x')\in P_{\Phi,q}^{\theta,C}:
Q_{q,j}^{\theta,C}(x)\neq Q_{q,j}^{\theta,C}(x')
\right\}.
\]

\begin{proposition}[Pair-cover characterization]
\label{prop:pair-cover}
A receiver family $K\subseteq I_{q+1}$ realizes $\Phi$ if and only if
\begin{equation}
P_{\Phi,q}^{\theta,C}
\subseteq
\bigcup_{j\in K}E_{q,j}^{\theta,C}.
\label{eq:pair-cover}
\end{equation}
Consequently,
\begin{equation}
d_{\Phi,q}^{\theta,C}
=
\min\left\{
|K|:
K\subseteq I_{q+1},\ 
P_{\Phi,q}^{\theta,C}
\subseteq
\bigcup_{j\in K}E_{q,j}^{\theta,C}
\right\},
\label{eq:pair-cover-order}
\end{equation}
where the minimum is $+\infty$ if no such $K$ exists. The empty family realizes exactly when $\Phi_q^{\theta,C}$ is constant.
\end{proposition}

\begin{proof}
By \eqref{eq:family-realization}, $K$ fails to realize $\Phi$ exactly when there is a pair $(x,x')$ distinguished by $\Phi_q^{\theta,C}$ that lies in every receiver kernel $D_{q,j}^{\theta,C}$ for $j\in K$. Equivalently, $K$ realizes $\Phi$ exactly when every pair in $P_{\Phi,q}^{\theta,C}$ is separated by at least one receiver in $K$, which is \eqref{eq:pair-cover}. The formula for $d_{\Phi,q}^{\theta,C}$ follows from \eqref{eq:realization-order}.
\end{proof}

Thus the localization order is the minimum number of declared receivers needed to cover all state-pair distinctions demanded by the phenomenon.

The count $d_{\Phi,q}^{\theta,C}$ is useful but incomplete. A circuit claim normally identifies components, not only their number. We therefore also care about the inclusion-minimal elements of $\Rcal_{\Phi,q}^{\theta,C}$. The next section shows that both the minimum count and the identity of these minimal families depend on support.

\section{Localization under support restriction}
\label{sec:support-restriction}

We now compare localization claims made on nested regimes while holding the network, receiver decomposition, and ambient phenomenon fixed.

Fix an ambient source regime $C_\ast\subseteq X_0$ and a phenomenon profile
\[
\Phi_r^{\ast}:S_r^{\theta,C_\ast}\to Y_\Phi.
\]
For each nonempty $C\subseteq C_\ast$, restrict $\Phi_r^{\ast}$ to $S_r^{\theta,C}$ and pull it back to earlier cuts. We write the resulting profile as $\Phi_q^{\theta,C}$.

If $C\subseteq C'\subseteq C_\ast$, then
\[
S_q^{\theta,C}\subseteq S_q^{\theta,C'}
\]
and both the receiver and phenomenon equivalence relations restrict to the smaller domain:
\begin{align}
D_{q,j}^{\theta,C}
&=
D_{q,j}^{\theta,C'}
\cap
\bigl(S_q^{\theta,C}\times S_q^{\theta,C}\bigr),
\label{eq:receiver-kernel-restriction}\\
D_{\Phi,q}^{\theta,C}
&=
D_{\Phi,q}^{\theta,C'}
\cap
\bigl(S_q^{\theta,C}\times S_q^{\theta,C}\bigr).
\label{eq:phenomenon-kernel-restriction}
\end{align}
The network has not acquired a new factorization. The criterion is being tested on fewer state pairs.

\begin{theorem}[Realization is preserved by support restriction]
\label{thm:realization-monotonicity}
Let $C\subseteq C'\subseteq C_\ast$. If a receiver family $K$ realizes $\Phi$ on $C'$, then $K$ realizes $\Phi$ on $C$. Equivalently,
\begin{equation}
\boxed{
\Rcal_{\Phi,q}^{\theta,C'}
\subseteq
\Rcal_{\Phi,q}^{\theta,C}.
}
\label{eq:realization-family-monotonicity}
\end{equation}
\end{theorem}

\begin{proof}
Assume $K\in\Rcal_{\Phi,q}^{\theta,C'}$. Restrict
\[
\bigcap_{j\in K}D_{q,j}^{\theta,C'}
\subseteq
D_{\Phi,q}^{\theta,C'}
\]
to $S_q^{\theta,C}\times S_q^{\theta,C}$ and use \eqref{eq:receiver-kernel-restriction}--\eqref{eq:phenomenon-kernel-restriction}.
\end{proof}

\begin{corollary}[Narrower support cannot increase the localization order]
\label{cor:order-monotonicity}
Under the hypotheses of \cref{thm:realization-monotonicity},
\begin{equation}
\boxed{
 d_{\Phi,q}^{\theta,C}
\le
 d_{\Phi,q}^{\theta,C'}.
}
\label{eq:order-monotonicity}
\end{equation}
\end{corollary}

The theorem gives a one-sided comparison. A broad regime demands that the proposed receiver family preserve the phenomenon across more possible internal distinctions. Restricting the regime removes some of those comparisons, so a family that failed globally may become sufficient locally. The reverse inference is invalid without additional information.

The continuous example in \cref{sec:supported-process} is already strict. For the linear readout $\Phi(x_1,x_2)=x_1+x_2$ with coordinate accesses, the full domain $\mathbb R^2$ requires both coordinates, while the diagonal $D=\{(t,t)\}$ has localization order one. The phenomenon remains nonconstant. The decrease is caused entirely by the relation $x_1=x_2$ on the restricted support.

\subsection{Minimal identity can change without changing the count}

Support also changes which receiver families are minimal. Consider the ambient phenomenon
\[
\Phi(x_1,x_2)=x_1
\]
on $\mathbb R^2$. The first coordinate is the unique one-coordinate realization on the full domain. On the diagonal $D$, either coordinate realizes the same phenomenon because $x_1=x_2$. The minimum count remains one, but the set of minimal realizations changes from
\[
\bigl\{\{1\}\bigr\}
\quad\text{to}\quad
\bigl\{\{1\},\{2\}\bigr\}.
\]
Circuit size therefore carries less information than circuit identity. Agreement in minimum size does not imply agreement in the components that realize the phenomenon, and disagreement in identity does not by itself imply that the ambient computation changed.

At the other extreme, a dependency can disappear completely when the phenomenon becomes constant on a restricted support. Taking parity itself as the phenomenon, the example from \cref{sec:support-dependency} gives such a case. These are exact changes in the feasible localization set, not failures of an extraction algorithm.

\subsection{Natural and intervention supports}

\Cref{thm:realization-monotonicity} is stated for source-generated regimes because those supports come with compatible continuations across cuts. At a fixed cut, its proof uses only restriction of the receiver and phenomenon relations. For arbitrary analysis supports $S\subseteq S'\subseteq X_q$ in the sense of \cref{def:analysis-support}, the same implication holds whenever the same ambient access and phenomenon are restricted to both domains.

Intervention protocols make this distinction concrete. Let $S_q^{\mathrm{nat}}$ be an analysis support generated by natural inputs, and let $S_q^{\mathrm{int}}$ contain states admitted by a patching, ablation, or other intervention protocol. The second set need not be the prefix image of any source regime. The statements
\[
Q_{q,j}^\theta|_{S_q^{\mathrm{nat}}}
\quad\text{realizes }\Phi|_{S_q^{\mathrm{nat}}}
\]
and
\[
Q_{q,j}^\theta|_{S_q^{\mathrm{int}}}
\quad\text{realizes }\Phi|_{S_q^{\mathrm{int}}}
\]
are therefore claims on different restricted domains.

If one support contains the other, exact realization transports from the larger domain to the smaller one. If the domains are incomparable, support restriction alone gives no implication in either direction. Intervention-induced hidden states can leave the natural representation distribution and can recruit pathways not active on ordinary inputs \citep{makelov2024subspace,grant2026divergent}. In the present formalism, those intervention states belong to the domain on which the localization claim is stated.

\section{Comparing localization across contexts}
\label{sec:context-comparison}

Support restriction gives an exact comparison when one regime is contained in another. Cross-context studies are often harder: two prompt classes or intervention regimes may overlap only partially, or not at all at the input level. The relevant overlap for a mechanistic claim can nevertheless occur later in the network.

\subsection{Input overlap can understate internal overlap}

Fix an ambient source regime $C_\ast$. For $C,D\subseteq C_\ast$, the internal supports at cut $q$ are
\[
S_q^{\theta,C}=A_q^\theta(C),
\qquad
S_q^{\theta,D}=A_q^\theta(D).
\]
Their shared source inputs generate only
\[
S_q^{\theta,C\cap D}=A_q^\theta(C\cap D).
\]
There can be additional shared internal states because the prefix $A_q^\theta$ need not be injective.

\begin{proposition}[Source intersection and internal overlap]
\label{prop:source-overlap}
For $C,D\subseteq C_\ast$ and every cut $q$,
\begin{equation}
S_q^{\theta,C\cap D}
\subseteq
S_q^{\theta,C}\cap S_q^{\theta,D}.
\label{eq:source-overlap-inclusion}
\end{equation}
If $A_q^\theta|_{C_\ast}$ is injective, equality holds for every $C,D\subseteq C_\ast$. If it is not injective, there exist $C,D\subseteq C_\ast$ for which the inclusion is strict.
\end{proposition}

\begin{proof}
The inclusion is the general relation $f(C\cap D)\subseteq f(C)\cap f(D)$ applied to $f=A_q^\theta$. If $A_q^\theta|_{C_\ast}$ is injective and $z$ lies in both images, then $z=A_q^\theta(x)=A_q^\theta(x')$ for some $x\in C$ and $x'\in D$. Injectivity gives $x=x'\in C\cap D$, so $z\in S_q^{\theta,C\cap D}$.

If the restriction is noninjective, choose distinct $x,x'\in C_\ast$ with $A_q^\theta(x)=A_q^\theta(x')=z$. Taking $C=\{x\}$ and $D=\{x'\}$ gives $C\cap D=\varnothing$ while $z\in S_q^{\theta,C}\cap S_q^{\theta,D}$.
\end{proof}

Thus two disjoint prompt classes can still pass through a common internal state. Comparing explanations only on the image of their shared input labels can miss part of the domain on which the internal claims genuinely overlap.

\subsection{Cross-support compatibility and union-level realization}

A localization can work on two regimes separately and fail when the regimes are combined. The obstruction can be stated exactly. Let $S_1,S_2\subseteq X$ be analysis supports, let
\[
L:S_1\cup S_2\twoheadrightarrow U
\]
be a fixed access corestricted to its realized image on the union, and let
\[
\Phi:S_1\cup S_2\to Y_\Phi
\]
be a fixed phenomenon. For $i\in\{1,2\}$, corestrict $L|_{S_i}$ to $L(S_i)$ when evaluating realization on $S_i$.

\begin{proposition}[Cross-support compatibility criterion]
\label{prop:cross-support-compatibility}
Assume $L|_{S_1}$ realizes $\Phi|_{S_1}$ and $L|_{S_2}$ realizes $\Phi|_{S_2}$. Then $L$ realizes $\Phi$ on $S_1\cup S_2$ if and only if
\begin{equation}
x\in S_1,\ y\in S_2,\ L(x)=L(y)
\quad\Longrightarrow\quad
\Phi(x)=\Phi(y).
\label{eq:cross-support-compatibility}
\end{equation}
\end{proposition}

\begin{proof}
By local realization, every pair in $S_1\times S_1$ or $S_2\times S_2$ that lies in $\Ker L$ already lies in $\Ker\Phi$. The only access-fiber pairs on $S_1\cup S_2$ not covered by those two local conditions are cross-support pairs. Hence
\[
\Ker(L|_{S_1\cup S_2})
\subseteq
\Ker(\Phi|_{S_1\cup S_2})
\]
holds exactly when every $x\in S_1$ and $y\in S_2$ with $L(x)=L(y)$ also satisfy $\Phi(x)=\Phi(y)$. The claim follows from \cref{thm:realization-criterion}.
\end{proof}

The condition checks the new comparisons introduced by the union. Local realization handles pairs internal to each support; union-level realization adds the requirement that access collisions across the two supports remain phenomenon-consistent.

\begin{example}[Local realization need not extend to a union]
\label{ex:local-not-global}
Let
\[
Q(a,b)=a,
\qquad
S_1=\{(0,0),(1,1)\},
\qquad
S_2=\{(0,0),(2,1)\},
\]
and take the access $L=\pi_{\{2\}}$. On $S_1$, the second coordinate determines $Q$: $b=0$ gives $Q=0$ and $b=1$ gives $Q=1$. On $S_2$, it again determines $Q$: $b=0$ gives $Q=0$ and $b=1$ gives $Q=2$. Thus $L$ realizes $\Phi=Q$ on each support, and trivially on their intersection.

On the union,
\[
S_1\cup S_2
=
\{(0,0),(1,1),(2,1)\},
\]
the cross-support pair $(1,1)\in S_1$ and $(2,1)\in S_2$ has
\[
L(1,1)=L(2,1)=1,
\qquad
Q(1,1)=1\neq 2=Q(2,1).
\]
The compatibility condition \eqref{eq:cross-support-compatibility} fails, so $L$ does not realize $Q$ on $S_1\cup S_2$.
\end{example}

For cross-context mechanistic analysis, the two checks are therefore separate. Internal overlap should be computed at the cut where the localization is stated, since source overlap can be too small. A localization that is valid separately on two supports extends to their union exactly when the cross-support access collisions satisfy \eqref{eq:cross-support-compatibility}.

\section{Discussion}
\label{sec:discussion}

\subsection{Typing localization claims}

A statement that behavior $\Phi$ is localized to receiver family $K$ has several arguments. For a source-generated support, the exact claim is
\[
K\in\Rcal_{\Phi,q}^{\theta,C}.
\]
It specifies the ambient network $\theta$, the represented cut $q$, the source regime $C$ that generates the supported domain, the phenomenon $\Phi$, and the declared receiver decomposition from which $K$ is drawn. At a fixed cut, the same definition applies on an arbitrary analysis support $S\subseteq X_q$ by restricting the access maps and phenomenon to $S$.

The domain enters the truth conditions of the claim. \Cref{thm:realization-monotonicity} gives a directional comparison when two supports are nested. \Cref{sec:context-comparison} shows why nonnested contexts require separate internal-overlap and cross-support compatibility checks. These statements distinguish a change in the ambient computation from a change in the state pairs that the localization must separate.

\subsection{Prompt-local analyses and counterfactual domains}

A singleton source regime in a deterministic network generates a singleton support at every reachable cut. Every exact sufficiency question on that literal support is then trivial: all maps are constant on a one-point domain. Prompt-local attribution, ablation, corruption, and patching methods become nontrivial by introducing additional states against which the prompt is compared.

Those baseline, corrupted, patched, or counterfactual states are therefore part of the effective domain of the mechanistic claim. Two analyses can begin from the same prompt and parameter setting while testing localization on different analysis supports. The distinction is formal rather than terminological: different supports contain different state pairs, and the realization criterion is evaluated on those pairs.

\subsection{Scope and limitations}

The theory is exact and set-theoretic. A probabilistic extension would need to distinguish exact support from high-probability regions and distribution-weighted approximation. Approximate realization would replace exact kernel inclusion with a quantitative condition on variation of the phenomenon within access fibers.

The access decomposition is declared by the analysis. Neuron-, head-, edge-, feature-, and mediator-level studies expose different internal states, and the present results are conditional on that choice.

The network is represented as a finite sequence of deterministic state cuts. Recurrent systems, stochastic computation, adaptive tool use, and training-time dynamics require a more general process index.

Exact interface realization is also weaker than causal necessity, counterfactual control, algorithmic abstraction, or equivalence between different presentation languages. Those questions require additional structure after the supported domain, phenomenon, and internal access have been specified.

\section{Conclusion}

A fixed neural network specifies ambient maps, but a mechanistic-localization claim is evaluated on a selected domain of internal states. The support belongs in the mathematical type of the claim together with the phenomenon and declared internal access.

With those arguments explicit, exact realization is a factorization problem on the supported domain. Restricting support removes state pairs from the comparison and can therefore enlarge the feasible localization family, lower its minimum receiver count, or change which families are minimal without changing the weights. Across nonnested contexts, source overlap can miss shared internal states, while separate local realizations extend to the union exactly when cross-support access collisions are phenomenon-consistent.

Mechanistic localizations should therefore be compared together with the domains on which they are asserted. Reporting only the network and circuit identity can conflate changes in ambient computation with changes in which internal distinctions the analysis requires the circuit to preserve.


\bibliography{references}
\bibliographystyle{tmlr}

\end{document}
